\documentclass[11pt]{amsart}
\usepackage{calc,amssymb,amsmath,ulem,stmaryrd,fullpage}
\usepackage{enumerate}
\usepackage{alltt}
\RequirePackage[dvipsnames,usenames]{color}
\let\mffont=\sf
\normalem
%\input{kmacros3.sty}
\input{xy}
\xyoption{all}
\usepackage{hyperref}
\usepackage{graphicx}
\usepackage[all,cmtip]{xy}
\usepackage{verbatim}
\usepackage[left=0.95in,top=0.95in,right=0.95in,bottom=0.95in]{geometry}
\newcommand{\tauCohomology}{T}
\newcommand{\FCohomology}{S}
\usepackage{mathrsfs}


\theoremstyle{theorem}
%\newtheorem*{mainthm*}{Main Theorem}

\newcommand{\note}[1]{\marginpar{\sffamily\tiny #1}}

\def\todo#1{\textcolor{Mahogany}%
{\sffamily\footnotesize\newline{\color{Mahogany}\fbox{\parbox{\textwidth-15pt}{#1}}}\newline}}

\renewcommand{\O}{\mathcal O}
\newcommand{\ie}{{\itshape i.e.} }
\newcommand{\roundup}[1]{\lceil #1 \rceil}
\newcommand{\rounddown}[1]{\lfloor #1 \rfloor}
\renewcommand{\to}[1][]{\xrightarrow{\ #1\ }}
\newcommand{\onto}[1][]{\protect{\xrightarrow{\ #1\ }\hspace{-0.8em}\rightarrow}}
\newcommand{\into}[1][]{\lhook \joinrel \xrightarrow{\ #1\ }}
%\newcommand{DBDual}[1]{{\underline{\omega}_{#1}}}
\newcommand{\DBDual}[1]{{{\uuline \omega}{}^{\mydot}_{#1}}}
\newcommand{\Tt}{{\mathfrak{T}}}
%\newcommand{\bn}{{\mathfrak{n}} }
\newcommand{\sol}[1]{ \vskip 6pt{\color{blue} {\bf Solution:} #1} \vskip 6pt}

\newcommand{\bR}{{\mathbb{R}}}
\newcommand{\va}{{\vec{a}}}
\newcommand{\vb}{{\vec{b}}}
\newcommand{\vzero}{{\vec{0}}}
\newcommand{\vi}{{\vec{i}}}
\newcommand{\vj}{{\vec{j}}}
\newcommand{\vk}{{\vec{k}}}
\newcommand{\vh}{{\vec{h}}}
\newcommand{\vr}{{\vec{r}}}
\newcommand{\vu}{{\vec{u}}}
\newcommand{\vv}{{\vec{v}}}
\newcommand{\vw}{{\vec{w}}}
\newcommand{\vx}{{\vec{x}}}
\newcommand{\vy}{{\vec{y}}}
\newcommand{\vz}{{\vec{z}}}
\newcommand{\vT}{{\vec{T}}}
\newcommand{\vN}{{\vec{N}}}
\newcommand{\vF}{{\vec{F}}}
\newcommand{\vG}{{\vec{G}}}
\newcommand{\bF}{\mathbb{F}}
\newcommand{\bQ}{\mathbb{Q}}
\newcommand{\bZ}{\mathbb{Z}}
\newcommand{\bA}{\mathbb{A}}
\newcommand{\cF}{\mathcal{F}}
\newcommand{\cG}{\mathcal{G}}
\newcommand{\cH}{\mathcal{H}}
\newcommand{\image}{\mathrm{im}\,}
\newcommand{\coker}{\mathrm{coker}\,}
\newcommand{\Hom}{\mathrm{Hom}\,}
\newcommand{\Supp}{\mathrm{Supp}\,}
\newcommand{\sH}{\mathscr{H}\,}

\begin{document}
\title{Worksheet \#1 -- Math 6140\\Spring 2019}
\author{Due Friday, January 18th}

\maketitle
You may work in groups of up to 3.  Only one worksheet needs to be turned in per group.
\vskip 9pt
\noindent
{\bf 1.}  Suppose that $\phi : \cF \to \cG$ is a map of sheaves (on $X$).  Determine which of the following presheaves are actually sheaves.  If they are, prove it.  If they are not, give an example.
\vskip 3pt
{\bf (a)}  the presheaf $\cH$ such that $\cH(U) = \ker\big(\cF(U) \to \cG(U)\big)$.
\vskip 120pt
{\bf (b)}  the presheaf $\cH(U) = \image\big(\cF(U) \to \cG(U)\big)$.
\vskip 120pt
{\bf (c)}  the presheaf $\cH(U) = \coker\big(\cF(U) \to \cG(U)\big) = \cG(U)/\cF(U)$.
\vskip 120pt
\vfill
The sheafification of (a),(b),(c) (if sheafification is necessary) are denoted $\ker \phi$, $\image \phi$ and $\coker \phi$ respectively.
\newpage
\noindent
{\bf 2.}  Suppose that $0 \to \cF \to \cG \to \cH \to 0$ is a short exact sequence of sheaves on $X$ (meaning it is short exact on the stalks).  Prove carefully that for every open $U \subseteq X$, the sequence:
\[
0 \to \cF(U) \to \cG(U) \to \cH(U)
\]
is exact.
\vskip 160pt
\noindent
{\bf 3.}  A sheaf $\cF$ is \emph{flasque} if for every $V \subseteq U$ (open sets in $X$), the restriction map $\cF(U) \to \cF(V)$ is surjective.  With the notation as in {\bf 2.}, show that if $\cF$ is {flasque}, then
\[
0 \to \cF(U) \to \cG(U) \to \cH(U) \to 0
\]
is exact.
\vskip 160pt
\noindent
{\bf 4.}  Suppose that $f : X \to Y$ is a map of topological spaces and $0 \to \cF \to \cG \to  \cH \to 0$ is a short exact sequence of sheaves on $X$.  Prove that
\[
0 \to f_* \cF \to f_* \cG \to f_* \cH
\]
is exact on $Y$.  Further give an example to show that $f_* \cG \to f_* \cH$ is not necessarily surjective.
\newpage
\noindent
{\bf 5.}  Suppose that $\phi : X \to Y$ is a regular map of affine varieties.  Suppose that $M$ is a $k[X]$ modules and $\cF = \widetilde{M}$.  Note that $M$ can also be viewed as a $k[Y]$ module by the map $\phi^* : k[Y] \to k[X]$.  We denote this module by $M_{k[Y]}$ and let $\cG = \widetilde{M_{k[Y]}}$.  Show that $\phi_* \cF \cong \cG$.
\vskip 200pt
\noindent
If $\phi : X \to Y$ is a continuous map of topological spaces and $\cG$ is a sheaf on $Y$, we define $\phi^{-1}(\cG)$ to be the sheafification of the presheaf $\cH$ where for any $U \subseteq X$, we form the direct limit:
\[
\cH(U) = \lim_{V \supseteq f(U)} \cG(V).
\]
\vskip 12pt
\noindent
{\bf 6.}  Suppose that $Z \subseteq X$ is a subset of a topological space and that $\cF$ is a sheaf on $X$.  Let $i : Z \to X$ be the inclusion and show that for each point $z \in Z$,
\[
(i^{-1}\cF)_z = \cF_z.
\]
In the case that $Z$ is a closed subvariety of a variety $X$, how do $i^{-1} \O_X$ and $\O_Z$ compare?
\newpage
\noindent
{\bf 7.}  Suppose that $\phi : X \to Y$ is a continuous map of topological spaces.  Show that there is a natural\footnote{Actually, please don't show the \emph{natural} part, but if you don't know what \emph{natural} means, ask.} bijection of sets for any sheaves $\cG$ on $Y$ and $\cF$ on $X$.
\[
\Hom_X(\phi^{-1} \cG, \cF) \cong \Hom_Y(\cG, \phi_* \cF).
\]
\vskip 200pt
First note that if $\cF$ is a sheaf, we define the \emph{support of $\cF$}, denoted $\Supp \cF$, to be $\{x \in X \;|\; \cF_x \neq 0\}$.  If $s \in \cF(U)$, we define $\Supp(s) = \{x \in U\;|\; s_x \neq 0\}$.  It is not difficult to show that $\Supp(s)$ is always closed but that $\Supp \cF$ is not always closed.
\vskip 12pt
\noindent
{\bf 8.}  Suppose that $Z \subseteq X$ is a closed subset and let $\cF$ be a sheaf on $X$.  Consider the presheaf $\sH^0_Z(\cF)$ defined as
\[
\sH^0_Z(\cF)(U) = \{s \in \cF(U)\;|\; \Supp s \subseteq Z \}.
\]
Show that $\sH^0Z(\cF)$ is a sheaf and that if $i : (X \setminus Z) \to X$ is the inclusion, there exists an exact sequence:
\[
0 \to \sH^0_Z(\cF) \to \cF \to i_* (\cF|_{X \setminus Z}).
\]
\end{document}

